\section{回归、加权与双重稳健估计各自依赖什么}
\label{sec:13-Machine-Learning/12-Causal-Inference/05}

\hyperref[sec:13-Machine-Learning/12-Causal-Inference/04]{上一节}用潜在结果把 ATE 定义清楚了，也给出了识别的条件。但识别只是回答``无限总体中的 ATE 可由观测分布中的哪些量表达''，它不负责有限样本怎么办。本节处理估计问题：已知

\[
(C_i,X_i,Y_i),\qquad i=1,\ldots,n,
\]

承认一致性、条件可忽略性和重叠条件已经成立，怎样从 \(n\) 条记录算出 ATE 的合理估计。

三层漏斗不能颠倒：第一层是因果假设是否成立；第二层是识别——目标能否写成观测分布的泛函；第三层才是估计。一个估计量方差很小，只说明它稳定地估计了某个量，不说明那个量就是 ATE。

\subsection{Outcome regression：先建模条件结果}

定义

\[
m_c(x)=\E[Y\mid C=c,X=x],
\]

即每种处理水平下的条件期望结果。拟合 \(\widehat m_1,\widehat m_0\) 后，用全体样本的协变量分布做标准化：

\[
\widehat\tau_{\mathrm{OR}}
=
\frac1n\sum_{i=1}^{n}
\left[
\widehat m_1(C_i)-\widehat m_0(C_i)
\right].
\]

思路很自然：为每个个体分别预测它在两种处理下的平均结果，再取总体平均差。这个估计量把全部赌注押在 outcome model 正确上。如果 \(m_x(c)\) 被正确指定，它可以有效利用结果与协变量之间的结构信息，方差通常比不加建模的方法小。但如果模型在稀少处理区域外推错误，估计就会系统偏离——而真正的数据里，处理组与对照组的协变量分布往往不重叠在同一区域，外推正是常态而非例外。

有一个常见的隐蔽错误：把 \(C\) 和 \(X\) 一起塞进线性回归，然后直接读 \(X\) 的系数当 ATE。只有在无交互、效应恒定且模型线性的强设定下这个系数才等于 ATE。当允许交互或非线性后，应显式计算标准化差而不是期待某个回归系数自动替你完成总体加权。

\subsection{Propensity score：描述谁会接受处理}

定义倾向得分

\[
e(c)=P(X=1\mid C=c).
\]

在条件可忽略性下，标量 \(e(C)\) 足以平衡 \(C\) 的处理组分布。逆概率加权的估计量为

\[
\widehat\tau_{\mathrm{IPW}}
=
\frac1n\sum_{i=1}^{n}
\left[
\frac{X_iY_i}{\widehat e(C_i)}
-
\frac{(1-X_i)Y_i}{1-\widehat e(C_i)}
\right].
\]

直觉是：某类人中处理很罕见，那么观察到的每个处理者就代表更多同类个体，给它更大的权重；对照罕见时同理。IPW 试图构造一个伪总体，其中处理分配与已测协变量无关。

Propensity model 追求的不是把处理---对照分得越准越好。如果模型几乎完美区分类别，往往意味着重叠很差——有些协变量组合几乎只出现一种处理——此时权重 \(1/e(C)\) 或 \(1/(1-e(C))\) 大到不可用。应当检查加权前后的协变量平衡、权重分布和有効样本量，而不是只报告倾向模型的分类 AUC。

\subsection{AIPW：把两个 nuisance model 接在一起}

令 \(\widehat m_x(c)\) 和 \(\widehat e(c)\) 分别估计结果模型与倾向模型。增广逆概率加权估计量 (AIPW) 的形式为

\[
\begin{aligned}
\widehat\tau_{\mathrm{AIPW}}
=\frac1n\sum_{i=1}^{n}
\Bigg[&
\widehat m_1(C_i)-\widehat m_0(C_i)\\
&+\frac{X_i\bigl(Y_i-\widehat m_1(C_i)\bigr)}
{\widehat e(C_i)}
-\frac{(1-X_i)\bigl(Y_i-\widehat m_0(C_i)\bigr)}
{1-\widehat e(C_i)}
\Bigg].
\end{aligned}
\]

第一行就是 outcome regression 的标准化预测。后两项是实际结果减去预测结果的残差，再按倾向加权。可以把这两项理解为对 outcome model 预测的修正：如果结果模型在某个区域预测偏了，实际观察到的 \(Y_i\) 会通过残差项把估计拉回来——只要倾向模型对该区域的权重合理。

在常见正则条件下，只要 outcome model 或 propensity model 其中一个一致，ATE 估计仍可一致。这就是``双重稳健''的含义：两个 nuisance model 允许错一个。它不意味着两个都可以随意错；也不修复隐藏混杂、错误定义的处理变量、或在重叠不存在的区域强行外推；更不能在 propensity 逼近 \(0\) 或 \(1\) 时控制残差项的高方差。

\subsection{机器学习拟合 nuisance 与推断之间的空隙}

树、boosting、kernel 或神经网络可以灵活估计 \(m_x\) 与 \(e\)，但在同一批数据上既拟合复杂 nuisance model 又计算修正项，会产生过拟合偏差。Cross-fitting 把样本分成若干折：在其他折上拟合 nuisance model，再对当前折计算预测与 influence-function 修正。这样每个观测的修正量来自未曾见过该观测的模型，缓解了过拟合渗入估计的问题。

但渐近正态和有效标准误还需要 nuisance 估计误差以足够快的速率衰减。Cross-fitting 减少了经验过程依赖，不等于任意黑箱模型都自动达到所需速率。报告时应区分四件事：识别假设成立的前提、nuisance 拟合策略及其调参、效应估计公式本身、以及不确定性的来源与计算方式。

\subsection{三个估计量各自敬畏不同的东西}

Outcome regression 敬畏结果模型的外推——模型错了，估计就错了。IPW 敬畏倾向得分的极端值——权重在重叠薄弱处决定一切。AIPW 同时敬畏两者，但给了第二条命：一个模型正确时能扛住另一个模型错设。

因此比较方法时不应只选置信区间最窄的那个。还要观察：协变量在加权前后是否真的平衡了；重叠区域覆盖了哪些样本；更换模型族或调参后效应估计是否稳定；权重截断规则改变了哪些个体的贡献；目标人群的定义改变是否翻转结论。稳定性分析的目的不是让多个方法相互投票选出一个得票最多的数字，而是定位：你报告的效应究竟由数据中的哪些样本、哪些假设和模型的哪个组件在支撑。
